% GENERATED FILE. Edit canonical lessons, then run npm run generate:books.
% canonical-source-hash: fnv1a64:8f4ad7c7114d89f6
% canonical-lessons: AR-W46-shin, AR-W46-sad, AR-W46-dal

\chapter{Letters From Your First Greetings}
\label{ch:ar-letters-from-greetings}
\hlchaptermodality{46}

\begin{chapteropening}
\textbf{By the end of this chapter:} \emph{I can write shin, sad and dal, and find each one in a greeting or thanks I have said since the first chapters.}

Complete the last lesson of chapter 46: i can write shin, sad and dal, and find each one in a greeting or thanks I have said since the first chapters.
\end{chapteropening}

\section[shīn]{\ar{ش} (shīn) — the letter that opens \ar{شكرا}}

\label{lesson:AR-W46-shin}



\begin{quote}
Before the new one: write \ar{ع} and \ar{ه} once each, and name them.

Now say \textbf{\ar{شكرا}} (\emph{shukran}), \textquotedblleft{}thank you\textquotedblright{}. You have read it for a long time. This lesson takes one letter out of it and puts it in your hand.
\end{quote}

\subsection*{Script you'll notice: \ar{ش}}
\textbf{\ar{ش}} — \emph{shīn}, the sound \emph{sh}.

It opens \textbf{\ar{شكرا}} (\emph{shukran}), \textquotedblleft{}thank you\textquotedblright{}, one of the first words in this book. You already write \textbf{\ar{س}} (\emph{sīn}): three small teeth and a bowl. \textbf{\ar{ش}} is the same body with \textbf{three dots above}. One skeleton, and the dots decide the sound.

\subsection*{Writing: \ar{ش}}
\hlblockfigure{\detokenize{figures/AR-W46-shin-filmstrip.pdf}}{How \ar{ش} is written, stroke by stroke}

\begin{itemize}
\raggedright
  \item \textbf{1.} start here: shape the three close teeth from right to left
  \item \textbf{2.} flow directly into the final bowl without lifting
  \item \textbf{3.} lift, then place the lower-left dot
  \item \textbf{4.} lift again, then place the lower-right dot
  \item \textbf{5.} lift a third time, then place the centered upper dot — and only now lift
\end{itemize}

\textbf{Pen lifts: 3.}

\begin{quote}
Stroke order is one attested teaching order; hands and teachers vary. Source: Introduction to Arabic: Egyptian Arabic for first-year students, Alphabet: \ar{س} \ar{ش} \ar{ص} \ar{ض}, independent Shiin at 00:00.7–00:03.0 (University of Oregon Libraries Open Textbook, 2023).
\end{quote}

\subsection*{Your turn}
\begin{itemize}
\raggedright
  \item \emph{Look:} at \textbf{\ar{شكرا}} and point at \ar{ش} inside it
  \item \emph{Trace it:} \ar{ش} three times, saying \emph{sh} as you finish each one
  \item \emph{Write it:} \ar{ه}, then \ar{ش}, from memory
  \item \emph{Read it:} \textbf{\ar{شكرا}} aloud once more
\end{itemize}

\begin{tcolorbox}[breakable,colback=teal!4,colframe=teal!35!black,title={Before you move on}]
Which letter is this — \ar{ش}? (\textbf{shīn}, \emph{sh}.) Which word did it come from? (\textbf{\ar{شكرا}}, \emph{shukran}, \textquotedblleft{}thank you\textquotedblright{}.)
\end{tcolorbox}
\section[ṣād]{\ar{ص} (ṣād) — the letter that opens \ar{صباح}}

\label{lesson:AR-W46-sad}



\begin{quote}
Before the new one: write \ar{ه} and \ar{ش} once each, and name them.

Now say \textbf{\ar{صباح} \ar{الخير}} (\emph{ṣabāḥ al-khayr}), \textquotedblleft{}good morning\textquotedblright{}. You have read it for a long time. This lesson takes one letter out of it and puts it in your hand.
\end{quote}

\subsection*{Script you'll notice: \ar{ص}}
\textbf{\ar{ص}} — \emph{ṣād}, a heavy \emph{s}.

It opens \textbf{\ar{صباح} \ar{الخير}} (\emph{ṣabāḥ al-khayr}), \textquotedblleft{}good morning\textquotedblright{}. Say \emph{ṣabāḥ} and feel the tongue press low and the back of the mouth widen: that heaviness is what the dot under the romanization (\emph{ṣ}) marks. The letter is a closed \textbf{oval} with a bowl trailing after it.

\subsection*{Writing: \ar{ص}}
\hlblockfigure{\detokenize{figures/AR-W46-sad-filmstrip.pdf}}{How \ar{ص} is written, stroke by stroke}

\begin{itemize}
\raggedright
  \item \textbf{1.} start here: close the oval clockwise from its lower-left junction
  \item \textbf{2.} turn left and rise into the short shoulder without lifting
  \item \textbf{3.} lift, restart at the baseline junction, and sweep through the trailing bowl — and only now lift
\end{itemize}

\textbf{Pen lifts: 1.}

\begin{quote}
Stroke order is one attested teaching order; hands and teachers vary. Source: Introduction to Arabic: Egyptian Arabic for first-year students, Alphabet: \ar{س} \ar{ش} \ar{ص} \ar{ض}, independent Saad at 00:01.1–00:03.3 (University of Oregon Libraries Open Textbook, 2023).
\end{quote}

\subsection*{Your turn}
\begin{itemize}
\raggedright
  \item \emph{Look:} at \textbf{\ar{صباح} \ar{الخير}} and point at \ar{ص} inside it
  \item \emph{Trace it:} \ar{ص} three times, saying \emph{ṣ} as you finish each one
  \item \emph{Write it:} \ar{ش}, then \ar{ص}, from memory
  \item \emph{Read it:} \textbf{\ar{صباح} \ar{الخير}} aloud once more
\end{itemize}

\begin{tcolorbox}[breakable,colback=teal!4,colframe=teal!35!black,title={Before you move on}]
Which letter is this — \ar{ص}? (\textbf{ṣād}, \emph{ṣ}.) Which word did it come from? (\textbf{\ar{صباح} \ar{الخير}}, \emph{ṣabāḥ al-khayr}, \textquotedblleft{}good morning\textquotedblright{}.)
\end{tcolorbox}
\section[dāl]{\ar{د} (dāl) — the letter that ends \ar{الحمد}}

\label{lesson:AR-W46-dal}



\begin{quote}
Before the new one: write \ar{ش} and \ar{ص} once each, and name them.

Now say \textbf{\ar{الحمد} \ar{لله}} (\emph{al-ḥamdu lillāh}), \textquotedblleft{}praise be to God\textquotedblright{}. You have read it for a long time. This lesson takes one letter out of it and puts it in your hand.
\end{quote}

\subsection*{Script you'll notice: \ar{د}}
\textbf{\ar{د}} — \emph{dāl}, the sound \emph{d}.

It closes \textbf{\ar{الحمد}} in \textbf{\ar{الحمد} \ar{لله}} (\emph{al-ḥamdu lillāh}), the everyday answer to \textquotedblleft{}how are you?\textquotedblright{}. \textbf{\ar{د}} is a \textbf{non-connector}, like \textbf{\ar{ا}} and \textbf{\ar{ر}}: it joins to the letter before it but never to the one after, so a small gap always follows it.

\subsection*{Writing: \ar{د}}
\hlblockfigure{\detokenize{figures/AR-W46-dal-filmstrip.pdf}}{How \ar{د} is written, stroke by stroke}

\begin{itemize}
\raggedright
  \item \textbf{1.} start here: begin at the upper tip and descend diagonally down and right through the curved shoulder
  \item \textbf{2.} turn left along the baseline without lifting — and only now lift
\end{itemize}

\textbf{Pen lifts: 0.} The pen never leaves the paper.

\begin{quote}
Stroke order is one attested teaching order; hands and teachers vary. Source: Introduction to Arabic: Egyptian Arabic for first-year students, Alphabet: \ar{د} \ar{ذ} \ar{ر}, independent Daal at 00:07.0–00:07.6 (University of Oregon Libraries Open Textbook, 2023).
\end{quote}

\subsection*{Your turn}
\begin{itemize}
\raggedright
  \item \emph{Look:} at \textbf{\ar{الحمد} \ar{لله}} and point at \ar{د} inside it
  \item \emph{Trace it:} \ar{د} three times, saying \emph{d} as you finish each one
  \item \emph{Write it:} \ar{ص}, then \ar{د}, from memory
  \item \emph{Read it:} \textbf{\ar{الحمد} \ar{لله}} aloud once more
\end{itemize}

\begin{tcolorbox}[breakable,colback=teal!4,colframe=teal!35!black,title={Before you move on}]
Which letter is this — \ar{د}? (\textbf{dāl}, \emph{d}.) Which word did it come from? (\textbf{\ar{الحمد} \ar{لله}}, \emph{al-ḥamdu lillāh}, \textquotedblleft{}praise be to God\textquotedblright{}.)
\end{tcolorbox}
